\documentclass[11pt,a4paper]{article}
\usepackage{gtu-el-yazisi}
\begin{document}
\notbilgi{Diferansiyel Denklemler}{Vize sonrası -- Cauchy--Euler, mertebe indirgeme, seri çözümleri, Laplace, konvolüsyon teoremi}{diferansiyel-not-vize-sonrasi-konvulasyon-teoremi.pdf}

\section{Cauchy-Euler denkleminin çözülmesi}
\sayfa{1}
\[ a_0x^ny^{(n)} + a_1x^{n-1}y^{(n-1)} + \cdots + a_{n-1}xy' + a_ny = f(x) \]
\[ \kutula{x=e^t} \qquad \kutula{t=\ln x} \]
Bu denklemin de karakteristik denklemi
\begin{multline*}
a_0k(k-1)\cdots(k-n+1) + a_1k(k-1)\cdots(k-n+2)\\
+ \cdots + a_{n-2}k(k-1) + a_{n-1}k + a_n = 0
\end{multline*}

\begin{ornek}
$x^2y'' + xy' + y = 4\sin(\ln x)$
\end{ornek}
\begin{enumerate}[label=\arabic*.]
\item Karakteristik denklem
\[ r(r-1) + r + 1 = 0 \;\rightarrow\; r^2+1 = 0 \]
\[ y''_t + y_t = 4\sin t \]
\item Homojen denklem
\[ y''_t + y_t = 0 \;\Rightarrow\; r^2+1 \;\Rightarrow\; r_1 = i,\ r_2 = -i \]
\[ y_h(t) = c_1\cos t + c_2\sin t \]
\end{enumerate}
Özel çözüm bulunur
\[ y_p(t) = t(A\cos t + B\sin t) \]
belirsiz katsayılar
\[ \Rightarrow\; y_p(t) = -2t\cos t \]

\sayfa{2}
O zaman
\[ y(t) = c_1\cos t + c_2\sin t - 2t\cos t \]
$t=\ln x$
\[ y(x) = c_1\cos(\ln x) + c_2\sin(\ln x) - 2\ln x\cos(\ln x)\,; \qquad c_1,c_2\in\mathbb{R} \]

\section{Mertebelerin indirgenmesi}
\sayfa{3}
\[ a_0(x)y^{(n)} + a_1(x)y^{(n-1)} + \cdots + a_{n-1}(x)y' + a_n(x)y = 0 \]
denklemi verilsin\\
$y_1 = y_1(x)$ çözümü bilinsin.

Bu durumda

``$y = u\cdot y_1$ \ ; \ $u=u(x)$'' dönüşümüyle

ve daha sonraki adımda

$u'=v$ dönüşümüyle denklemin mertebesi düşer. Böylece $(n-1)$. mertebe denklem elde edilir

\section{Hatırlatma}
\sayfa{4}
\[ \text{Taylor serisi} = \sum_{n=0}^{\infty} \frac{f^{(n)}(x_0)}{n!}\,(x-x_0)^n \]
$\ast$ Eğer $x=x_0$ noktasında $f(x)$ fonksiyonunun taylor serisi $x_0$ noktasının civarında yakınsaksa ve yakınsak olduğu fonksiyon $f(x)$'e eşitse, ``$f(x)$ fonksiyonu $x_0$ noktasında analitiktir'' denir.

\section{Serisel Çözümler}
\[ y'' + p(x)y' + q(x)y = 0 \]
2. mertebeden lineer ve homojen diferansiyel denklem verilsin.
\begin{itemize}[label=--]
\item Eğer, $p(x)$ ve $q(x)$ katsayıları bir $x_0$ noktasında \vurgu[pembe]{analitik} olurlarsa $x_0$ noktasına denklemin \vurgu[pembe]{adi noktası} denir.
\item Eğer, $p(x)$ ve $q(x)$ fonksiyonlarından birisi veya herikisi $x_0$ noktasında analitik değillerse, $x_0$ noktasına denklemin ``tekil (singüler) noktası'' denir.
\end{itemize}

\sayfa{5}
$\ast$ Eğer $x=x_0$ noktası denklemin adi noktası olursa, bu durumda denklemin her çözümü bir $|x-x_0|<r$ \ $(r>0)$ aralığında yakınsak olan
\[ \sum_{n=0}^{\infty} c_n(x-x_0)^n \quad \text{şeklinde} \]
fonksiyonlardır.

\subsection{Adi nokta civarında çözüm yolu}
\begin{enumerate}[label=\arabic*.]
\item Denklemi normal forma dönüştür yani
\[ y'' + p(x)y' + q(x)y = 0 \]
\item Bir adi nokta $x_0$ olarak seçilir.
\item $\displaystyle y = \sum_{n=0}^{\infty} c_n(x-x_0)^n$ şeklinde seçilir
\end{enumerate}
$c_n$ katsayılarını bulmak için türevleri alıp, denklemde (ilk denklem) yerine yazarız.\\
(Not $x_0\neq 0$ dan ise $x-x_0=t$ denilerek çözüme ulaşılır böylece daha kolay olur)

\sayfa{6}
\begin{enumerate}[label=\arabic*.,start=4]
\item katsayılar bulunduktan sonra serinin birkaç terimini yazarız
\item Başlangıç değer problemi ise başlangıç koşulları koyulup kalan katsayılar bulunur
\item $t=x-x_0$ koyulur ve $y(x)$ bulunur
\item genel çözüm\\
sırasıyla \quad $c_0=0,\ c_1=1$\\
sonra \quad $c_0=1$ ve $c_1=0$\\
vererek iki temel çözüm bulunur

ve çözüm bunların lineer kombinasyonudur
\end{enumerate}

\sayfa{7}
$\ast$ Eğer $x=x_0$ denklemin bir tekil noktasıysa ancak $(x-x_0)p(x)$ ve $(x-x_0)^2q(x)$ fonksiyonları $x=x_0$ noktasında analitik fonksiyonlar olursa $x=x_0$ noktasına denklemin ``regular tekil noktası'' denir.

\subsection{Forbenius yöntemi;}
Eğer $x=x_0$ noktası denklemin regular tekil noktası olursa bu nokta civarında denklemin
\[ y(x) = |x-x_0|^r \sum_{n=0}^{\infty} c_n(x-x_0)^n, \quad c_0\neq 0 \]
genelleştirilmiş kuvvet serisi şeklinde çözümü var oluyor.

$\ast$ Genelleştirilmiş kuvvet serisi yardımıyla çözüm yapılırken öncelikle
\[ \kutula{r(r-1) + p_0r + q_0} = 0 \]
denkleminden $r_1$ ve $r_2$ bulunur.
\[ p_0 = \lim_{x\to 0} x\,p(x) \quad\text{ve}\quad q_0 = \lim_{x\to 0} x^2q(x) \ \text{dir.} \]
Bu denkleme $x=0$ regüler tekil noktası için ``indis denklemi'' denir.

\sayfa{8}
Daha sonrasında aşağıdaki durumlar geçerli olur

\noindent 1) \underline{$r_1-r_2\neq 0,1,2,3,\cdots$ (negatif olmayan tamsayı değil)} ise
\begin{align*}
y_1(x) &= x^{r_1}\cdot\sum_{n=0}^{\infty} c_nx^n, \quad c_0\neq 0 \quad\text{ve}\\
y_2(x) &= x^{r_2}\cdot\sum_{n=0}^{\infty} c_n^{*}x^n, \quad c_0^{*}\neq 0
\end{align*}
2) \underline{$r_1-r_2 = 1,2,3,\cdots$ (doğal sayılar)} ise
\begin{align*}
y_1(x) &= x^{r_1}\cdot\sum_{n=0}^{\infty} c_nx^n, \quad c_0\neq 0\\
&\text{ve}\\
y_2(x) &= x^{r_2}\sum_{n=0}^{\infty} c_n^{*}x^n + c\,y_1(x)\ln x \qquad c_0^{*}\neq 0,\ \ c \text{ constant}
\end{align*}
3) \underline{$r_1-r_2=0$ ise}
\begin{align*}
y_1(x) &= x^{r_1}\cdot\sum_{n=0}^{\infty} c_nx^n, \quad c_0\neq 0 \quad\text{ve}\\
y_2(x) &= x^{r_1+1}\sum_{n=0}^{\infty} c_n^{*}x^n + y_1(x)\ln x
\end{align*}
(2) ve (3) durumlarında $y_2(x)$ çözümleri $y_1(x)$ çözümü bilindikten sonra mertebe düşürmekle elde edilir

\section{Laplace}
\sayfa{9}
Bir $f(x)$, $x>0$ fonksiyonu verilmiş olsun. Bu durumda,
\[ \kutula{F(s) = \int_0^{\infty} e^{-sx}f(x)\,dx} \]
integrali tanımlı ise bu $F$ fonksiyonuna $f$'in Laplace dönüşümü denir.
\[ F(s) = \mathcal{L}(f(x) \qquad\text{ve}\qquad \mathcal{L}^{-1}(F(s)) = f(x) \]
$\ast$ $f(x)$, $x>0$ fonksiyonu ve $a>0$ için
\[ \begin{cases} 0 & 0<x<a\\ f(x-a) & x>a \end{cases} \qquad \text{fonksiyonuna} \]
``geçiş fonksiyonu'' denir.

Birim basamak fonksiyonun
\[ u_a(x) = \begin{cases} 0, & 0<x<a\\ 1, & x>a \end{cases} \qquad \text{tanımına göre} \]

\sayfa{10}
geçiş fonksiyonu için
\[ u_a(x)\,f(x-a) = \begin{cases} 0, & 0<x<a\\ f(x-a), & x>a \end{cases} \quad \text{olur} \]
$\ast$ Bazı fonksiyonların Laplace dönüşümleri aşağıdaki gibidir
\begin{center}
\renewcommand\arraystretch{1.5}
\begin{tabular}{c|l}
$f(x),\ x>0$ & \quad $F(s)=\mathcal{L}(f(x))$ \\ \hline
$1$ & \quad $1/s$ \\
$x$ & \quad $1/s^2$ \\
$x^n$ & \quad $n!/s^{n+1}$ \\
$e^{ax}$ & \quad $1/s-a$ \\
$\sin(ax)$ & \quad $\dfrac{a}{s^2+a^2}$ \\
$\cos(ax)$ & \quad $\dfrac{s}{s^2+a^2}$ \\[2mm]
\multicolumn{1}{l|}{Özellikler} & \\
$e^{ax}f(x)$ & \quad $F(s-a)$ \\
$f^{(n)}(x)$ & \quad $s^nF(s) - s^{n-1}f(0)$ \\
 & \quad $-\,s^{n-2}f'(0) - \cdots$ \\
 & \quad $-\,f^{(n-1)}(0)$ \\
$x^nf(x)$ & \quad $(-1)^nF^{(n)}(s)$ \\
$u_a(x)f(x-a)$ & \quad $e^{-as}F(s),\ a>0$
\end{tabular}
\end{center}

\section{Konvulasyon teoremi}
\sayfa{11}
$\ast$ Eksponsiyonel dereceli ve her $0\le x\le b$ aralığında parçalı sürekli $f$ ve $g$ fonksiyonlarının konvolusyonu
\formulkutu{f(x) * g(x) = \int_0^{x} f(z)\,g(x-z)\,dz}
integrali ile belirlenmiş $f*g$ fonksiyonuna denir.

Ayrıca, \ $f*g = g*f$ dir

$\ast$ Eğer $f$ ve $g$ fonksiyonları $e^{at}$ exsponansiyonel dereceli ve her $0\le x\le b$ aralığında parçalı sürekli iki fonksiyon olursa, o halde $f*g$ konvolüsyonun Laplace dönüşümü
\[ \mathcal{L}[f*g] = \mathcal{L}[f]\cdot\mathcal{L}[g] \]
ve
\[ \mathcal{L}^{-1}[F(s)\cdot G(s)] = \mathcal{L}^{-1}[F(s)] * \mathcal{L}^{-1}[G(s)] \]

\end{document}
